\section{Conclusion}
\label{sec:conclusion}

We established the theoretical and empirical complexity profile of METIS
recursive bisection for congressional redistricting.

The key results are:

\begin{enumerate}
  \item \textbf{NP-hardness}: Balanced $k$-cut redistricting on planar
    graphs is NP-hard (Theorem~\ref{thm:nphard}), by reduction from
    Connected Planar Graph Bisection~\citep{dyer1985partitioning}.
    This correctly accounts for the contiguity requirement of the
    redistricting problem.  The hardness is for optimisation, not
    feasibility: a valid partition always exists for connected graphs.
  \item \textbf{Approximation landscape}: No formal approximation ratio
    is known for METIS's heuristic (Proposition~\ref{prop:approx}).
    The best known polynomial-time result ($O(\sqrt{\log n})$ via
    ARV~\citep{arora2009expander}) applies to Sparsest Cut, not to
    balanced connected $k$-partition.
    Empirically, METIS is within 3\% of the best of 10,000 seeds
    across redistricting-scale instances~\citep{b7paper}.
  \item \textbf{Runtime}: $O(\texttt{niter} \cdot n \log k)$ with
    $\texttt{niter} = 100$ fixed; empirically $O(n^{1.07 \pm 0.03})$
    on $k > 1$ states across three census years, confirming near-linear
    practical scaling.
    Single-district states are excluded from the scaling fit.
  \item \textbf{Space}: $O(n \log k)$, under 3\,MB for the largest
    state (California, $n = 8{,}057$, $k = 52$) at tract resolution.
\end{enumerate}

These results provide the complexity-theoretic foundation for the DIA's
algorithm specification.  NP-hardness establishes that exact optimisation
is intractable, making heuristic approximation not merely acceptable but
necessary.  The empirical 3\% gap from the best feasible seed provides
a practical bound on solution quality for legal and policy purposes.

\paragraph{Future work.}
Open problems include: (1) a tighter approximation ratio specific to the
redistricting graph class, (2) empirical scaling at block-level resolution
($n \approx 8.1\mathrm{M}$), and (3) complexity analysis of the
multi-constraint variant (VRA compliance + county preservation).
